\documentclass[10pt,a4paper]{amsart}
\usepackage{amsmath,amssymb,amsthm,graphicx}
% Retained preamble for the Stacks Project chapter morphisms.tex (Morphisms of Schemes).
% The chapter file pinned-original-source.tex carries only \input{preamble} on its first line, so no
% preamble text lives inside the chapter itself.  The Stacks Project uses one single shared preamble for
% every chapter; the snapshot stages it verbatim in sources/stacks-chow/preamble.tex.  The block below is
% copied mechanically, line for line, from that file's lines 14-18 (the Xy-pic setup) and lines 154-217
% (the native theorem declarations together with the mathematical macros).  Nothing was added, removed or
% redefined.  The wrapper class and the amsmath/amssymb/amsthm/graphicx base packages are supplied by the
% assembler, so the class/\externaldocument part of the shared preamble is not repeated here.
% For commutative diagrams we use Xy-pic
\usepackage[all]{xy}

% We use 2cell for 2-commutative diagrams.
\xyoption{2cell}
% Theorem environments.
%
\theoremstyle{plain}
\newtheorem{theorem}[subsection]{Theorem}
\newtheorem{proposition}[subsection]{Proposition}
\newtheorem{lemma}[subsection]{Lemma}

\theoremstyle{definition}
\newtheorem{definition}[subsection]{Definition}
\newtheorem{example}[subsection]{Example}
\newtheorem{exercise}[subsection]{Exercise}
\newtheorem{situation}[subsection]{Situation}

\theoremstyle{remark}
\newtheorem{remark}[subsection]{Remark}
\newtheorem{remarks}[subsection]{Remarks}

\numberwithin{equation}{subsection}

% Macros
%
\def\lim{\mathop{\mathrm{lim}}\nolimits}
\def\colim{\mathop{\mathrm{colim}}\nolimits}
\def\Spec{\mathop{\mathrm{Spec}}}
\def\Hom{\mathop{\mathrm{Hom}}\nolimits}
\def\Ext{\mathop{\mathrm{Ext}}\nolimits}
\def\SheafHom{\mathop{\mathcal{H}\!\mathit{om}}\nolimits}
\def\SheafExt{\mathop{\mathcal{E}\!\mathit{xt}}\nolimits}
\def\Sch{\mathit{Sch}}
\def\Mor{\mathop{\mathrm{Mor}}\nolimits}
\def\Ob{\mathop{\mathrm{Ob}}\nolimits}
\def\Sh{\mathop{\mathit{Sh}}\nolimits}
\def\NL{\mathop{N\!L}\nolimits}
\def\CH{\mathop{\mathrm{CH}}\nolimits}
\def\proetale{{pro\text{-}\acute{e}tale}}
\def\etale{{\acute{e}tale}}
\def\QCoh{\mathit{QCoh}}
\def\Ker{\mathop{\mathrm{Ker}}}
\def\Im{\mathop{\mathrm{Im}}}
\def\Coker{\mathop{\mathrm{Coker}}}
\def\Coim{\mathop{\mathrm{Coim}}}

% Boxtimes
%
\DeclareMathSymbol{\boxtimes}{\mathbin}{AMSa}{"02}

%
% Macros for moduli stacks/spaces
%
\def\QCohstack{\mathcal{QC}\!\mathit{oh}}
\def\Cohstack{\mathcal{C}\!\mathit{oh}}
\def\Spacesstack{\mathcal{S}\!\mathit{paces}}
\def\Quotfunctor{\mathrm{Quot}}
\def\Hilbfunctor{\mathrm{Hilb}}
\def\Curvesstack{\mathcal{C}\!\mathit{urves}}
\def\Polarizedstack{\mathcal{P}\!\mathit{olarized}}
\def\Complexesstack{\mathcal{C}\!\mathit{omplexes}}
% \Pic is the operator that assigns to X its picard group, usage \Pic(X)
% \Picardstack_{X/B} denotes the Picard stack of X over B
% \Picardfunctor_{X/B} denotes the Picard functor of X over B
\def\Pic{\mathop{\mathrm{Pic}}\nolimits}
\def\Picardstack{\mathcal{P}\!\mathit{ic}}
\def\Picardfunctor{\mathrm{Pic}}
\def\Deformationcategory{\mathcal{D}\!\mathit{ef}}

% --- appended 2026-10-12 (Chief): the project own definitions of the reference/slogan/history
% display environments, copied mechanically from the shared Stacks preamble lines 8-12
% (\usepackage{verbatim} plus the three \newenvironment lines, which the chapter preamble keeps
% verbatim from the project preamble). Nothing added, removed or redefined: these environments
% carry no mathematical content and the project itself typesets them as comments. Without these
% three lines a counted statement whose own text uses them fails with Environment slogan undefined.
\usepackage{verbatim}
\newenvironment{reference}{\comment}{\endcomment}
\newenvironment{slogan}{\comment}{\endcomment}
\newenvironment{history}{\comment}{\endcomment}
\usepackage{cleveref}
\crefname{originallemma}{original lemma}{original lemmas}
\Crefname{originallemma}{Original lemma}{Original lemmas}
\makeatletter
\newlabel{morphisms-lemma-characterize-normalization}{{035I}{source}{}{sourcepointer-1}{}}
\newlabel{morphisms-lemma-characterize-normalization@cref}{{[originallemma][0][]035I}{[1][0][]source}{}{}{}}
\makeatother
\crefname{originallemma}{original lemma}{original lemmas}
\Crefname{originallemma}{Original lemma}{Original lemmas}
\makeatletter
\newlabel{definition-normalization-X-in-Y}{{035H}{source}{}{sourcepointer-2}{}}
\newlabel{definition-normalization-X-in-Y@cref}{{[originallemma][0][]035H}{[1][0][]source}{}{}{}}
\makeatother
\crefname{originallemma}{original lemma}{original lemmas}
\Crefname{originallemma}{Original lemma}{Original lemmas}
\makeatletter
\newlabel{definition-integral}{{01WH}{source}{}{sourcepointer-3}{}}
\newlabel{definition-integral@cref}{{[originallemma][0][]01WH}{[1][0][]source}{}{}{}}
\makeatother
\crefname{originallemma}{original lemma}{original lemmas}
\Crefname{originallemma}{Original lemma}{Original lemmas}
\makeatletter
\newlabel{lemma-integral-closure}{{035F}{source}{}{sourcepointer-4}{}}
\newlabel{lemma-integral-closure@cref}{{[originallemma][0][]035F}{[1][0][]source}{}{}{}}
\makeatother
\crefname{originallemma}{original lemma}{original lemmas}
\Crefname{originallemma}{Original lemma}{Original lemmas}
\makeatletter
\newlabel{lemma-affine-equivalence-algebras}{{01SA}{source}{}{sourcepointer-5}{}}
\newlabel{lemma-affine-equivalence-algebras@cref}{{[originallemma][0][]01SA}{[1][0][]source}{}{}{}}
\makeatother
\crefname{originallemma}{original lemma}{original lemmas}
\Crefname{originallemma}{Original lemma}{Original lemmas}
\makeatletter
\newlabel{schemes-lemma-quasi-compact-permanence}{{03GI}{source}{}{sourcepointer-6}{}}
\newlabel{schemes-lemma-quasi-compact-permanence@cref}{{[originallemma][0][]03GI}{[1][0][]source}{}{}{}}
\makeatother
\crefname{originallemma}{original lemma}{original lemmas}
\Crefname{originallemma}{Original lemma}{Original lemmas}
\makeatletter
\newlabel{schemes-lemma-compose-after-separated}{{01KV}{source}{}{sourcepointer-7}{}}
\newlabel{schemes-lemma-compose-after-separated@cref}{{[originallemma][0][]01KV}{[1][0][]source}{}{}{}}
\makeatother
\crefname{originallemma}{original lemma}{original lemmas}
\Crefname{originallemma}{Original lemma}{Original lemmas}
\makeatletter
\newlabel{lemma-finite-permanence}{{035D}{source}{}{sourcepointer-8}{}}
\newlabel{lemma-finite-permanence@cref}{{[originallemma][0][]035D}{[1][0][]source}{}{}{}}
\makeatother
\title{Characterisation of the normalization by a universal property}
\author{The Stacks Project Authors}
\date{stacks-project:morphisms.tex:local-snapshot-20261012; GNU FDL 1.2 or later}
\begin{document}
\maketitle
\setcounter{subsection}{1}
\noindent\textit{Editorial note.} All delivered mathematical text is the author's own, copied from the pinned Stacks Project chapter file morphisms.tex; no statement or proof step was rewritten, reordered or removed.\par\medskip
\noindent\textit{Editorial note.} The chapter file pinned-original-source.tex carries only \textbackslash{}input\{preamble\} on its first line, so no preamble text lives inside the chapter. The retained file sources/stacks-morphisms/retained-preamble.tex is a mechanical, line-for-line copy of the project's single shared preamble as staged in the snapshot. Nothing was added, removed or redefined, and the excerpt serves only as preamble, never as a span coordinate.\par\medskip
\noindent\textit{Editorial note.} Both counted spans start at the \textbackslash{}begin line of their environment, so the statement and the author's proof are delivered whole; surrounding prose of the section is deliberately outside the counted spans and is carried as standard prerequisites instead.\par\medskip
\noindent\textit{Editorial note.} The citation fields (source\_id, source\_version, authors, source\_url, license.name, source\_title) are identical to those of every other card this lane has produced from this chapter, so the batch selector treats them as one source.\par\medskip
\noindent\textit{Editorial note.} Read in full under the Lead's 12:4x rule (screening hit re-read, pb 2239B). The author writes verbatim 'We omit the verification that the diagram commutes.' - a routine commutativity check for the constructed morphism h : X' -> Z; the construction of h and the subsequent argument are given in full.\par\medskip
\noindent\textit{Allowed standard prerequisites (editorial).}\begin{itemize}
\item Equivalence between quasi-coherent sheaves of algebras and affine morphisms. Lemma \textbackslash{}ref\{lemma-affine-equivalence-algebras\}, cited at proof lines 13969-13975. Provides the morphism h and its uniqueness.
\item Integral and birational morphisms are described affinely by integral extensions of coordinate rings. Lemma \textbackslash{}ref\{lemma-normalization\} and Lemma \textbackslash{}ref\{lemma-birational-integral\}, used at proof lines 13957-13968 and 13988-13999. Reduce the universal property to the affine case.
\end{itemize}\medskip
\ifdefined\hypertarget\hypertarget{sourcepointer-1}{}\fi
\noindent\textit{Source reference.} Original lemma 035I: focused-seeds/AGP/sources/stacks-snapshot-20261011/tags.txt, source lines 4129--4129. The lemma's own canonical tag, quoted so that readers can resolve the label inside the delivered text; it points at the counted statement itself and carries no proof content.\par\medskip
\ifdefined\hypertarget\hypertarget{sourcepointer-2}{}\fi
\noindent\textit{Source reference.} Original lemma 035H: focused-seeds/AGP/sources/stacks-snapshot-20261011/tags.txt, source lines 4128--4128. Navigation pointer to a statement the counted proof refers to but does not contain; nothing load-bearing is delegated to it.\par\medskip
\ifdefined\hypertarget\hypertarget{sourcepointer-3}{}\fi
\noindent\textit{Source reference.} Original lemma 01WH: focused-seeds/AGP/sources/stacks-snapshot-20261011/tags.txt, source lines 2482--2482. Navigation pointer to a statement the counted proof refers to but does not contain; nothing load-bearing is delegated to it.\par\medskip
\ifdefined\hypertarget\hypertarget{sourcepointer-4}{}\fi
\noindent\textit{Source reference.} Original lemma 035F: focused-seeds/AGP/sources/stacks-snapshot-20261011/tags.txt, source lines 4126--4126. Navigation pointer to a statement the counted proof refers to but does not contain; nothing load-bearing is delegated to it.\par\medskip
\ifdefined\hypertarget\hypertarget{sourcepointer-5}{}\fi
\noindent\textit{Source reference.} Original lemma 01SA: focused-seeds/AGP/sources/stacks-snapshot-20261011/tags.txt, source lines 2331--2331. Navigation pointer to a statement the counted proof refers to but does not contain; nothing load-bearing is delegated to it.\par\medskip
\ifdefined\hypertarget\hypertarget{sourcepointer-6}{}\fi
\noindent\textit{Source reference.} Original lemma 03GI: focused-seeds/AGP/sources/stacks-snapshot-20261011/tags.txt, source lines 4529--4529. Navigation pointer to a statement the counted proof refers to but does not contain; nothing load-bearing is delegated to it.\par\medskip
\ifdefined\hypertarget\hypertarget{sourcepointer-7}{}\fi
\noindent\textit{Source reference.} Original lemma 01KV: focused-seeds/AGP/sources/stacks-snapshot-20261011/tags.txt, source lines 2064--2064. Navigation pointer to a statement the counted proof refers to but does not contain; nothing load-bearing is delegated to it.\par\medskip
\ifdefined\hypertarget\hypertarget{sourcepointer-8}{}\fi
\noindent\textit{Source reference.} Original lemma 035D: focused-seeds/AGP/sources/stacks-snapshot-20261011/tags.txt, source lines 4124--4124. Navigation pointer to a statement the counted proof refers to but does not contain; nothing load-bearing is delegated to it.\par\medskip
\begin{lemma}
\label{lemma-characterize-normalization}
Let $f : Y \to X$ be a quasi-compact and quasi-separated morphism of schemes.
The factorization $f = \nu \circ f'$, where $\nu : X' \to X$ is the
normalization of $X$ in $Y$ is characterized by the following
two properties:
\begin{enumerate}
\item the morphism $\nu$ is integral, and
\item for any factorization $f = \pi \circ g$, with $\pi : Z \to X$
integral, there exists a commutative diagram
$$
\xymatrix{
Y \ar[d]_{f'} \ar[r]_g & Z \ar[d]^\pi \\
X' \ar[ru]^h \ar[r]^\nu & X
}
$$
for some unique morphism $h : X' \to Z$.
\end{enumerate}
Moreover, the morphism $f' : Y \to X'$ is dominant and in (2) the
morphism $h : X' \to Z$ is the normalization of $Z$ in $Y$.
\end{lemma}

\begin{proof}
Let $\mathcal{O}' \subset f_*\mathcal{O}_Y$ be the integral closure of
$\mathcal{O}_X$ as in Definition \ref{definition-normalization-X-in-Y}.
The morphism $\nu$ is integral by construction, which proves (1).
Assume given a factorization $f = \pi \circ g$ with $\pi : Z \to X$
integral as in (2). By Definition \ref{definition-integral}
$\pi$ is affine, and hence $Z$ is the relative
spectrum of a quasi-coherent sheaf of $\mathcal{O}_X$-algebras $\mathcal{B}$.
The morphism $g : Y \to Z$ corresponds to a map of $\mathcal{O}_X$-algebras
$\chi : \mathcal{B} \to f_*\mathcal{O}_Y$. Since $\mathcal{B}(U)$ is
integral over $\mathcal{O}_X(U)$ for every affine open $U \subset X$
(by Definition \ref{definition-integral})
we see from Lemma \ref{lemma-integral-closure}
that $\chi(\mathcal{B}) \subset \mathcal{O}'$.
By the functoriality of the relative spectrum
Lemma \ref{lemma-affine-equivalence-algebras}
this provides us with a unique morphism
$h : X' \to Z$. We omit the verification that the diagram commutes.

\medskip\noindent
It is clear that (1) and (2) characterize the factorization $f = \nu \circ f'$
since it characterizes it as an initial object in a category.

\medskip\noindent
From the universal property in (2) we see that $f'$ does not factor through
a proper closed subscheme of $X'$. Hence the scheme theoretic image of $f'$
is $X'$. Since $f'$ is quasi-compact (by Schemes, Lemma
\ref{schemes-lemma-quasi-compact-permanence} and the fact that $\nu$ is
separated as an affine morphism) we see that $f'(Y)$ is dense in $X'$.
Hence $f'$ is dominant.

\medskip\noindent
Observe that $g$ is quasi-compact and quasi-separated by
Schemes, Lemmas \ref{schemes-lemma-compose-after-separated} and
\ref{schemes-lemma-quasi-compact-permanence}. Thus the last statement
of the lemma makes sense.
The morphism $h$ in (2) is integral by Lemma \ref{lemma-finite-permanence}.
Given a factorization $g = \pi' \circ g'$ with $\pi' : Z' \to Z$
integral, we get a factorization $f = (\pi \circ \pi') \circ g'$ and
we get a morphism $h' : X' \to Z'$. Uniqueness implies that
$\pi' \circ h' = h$. Hence the characterization (1), (2) applies
to the morphism $h : X' \to Z$ which gives the final assertion of the lemma.
\end{proof}

\end{document}
